\subsection{Landscape of Modular Symmetric Flavor Models --- arXiv:2011.09154}

\textbf{Authors:} Keiya Ishiguro, Tatsuo Kobayashi, Hajime Otsuka

\paragraph{主な主張}
フラックスで固定される複素構造モジュライの分布は残存\(\mathbb{Z}_3\)対称性の固定点に集中し、CPを破る真空は相対的に少ない一方、その位相はaxio-dilatonの位相と相関することを示した。 \cite{Ishiguro:2020tmo}

\paragraph{新規性}
モジュラー・フレーバー模型で選ぶモジュライの真空値を、フラックス真空の統計から評価し、好まれる点とCP破れの分布を具体的に比較した。

\paragraph{位置付け}
フレーバー模型のモジュライ設定に対し、弦理論の安定化機構から統計的な指針を与える研究である。
